% === I. INTRODUCTION =========================================================
\section{Introduction}

Model-based controllers for autonomous racing derive their performance from the
fidelity of the vehicle model they carry~\cite{kabzan2019lbmpc}, and in the
lateral channel that fidelity is dominated by the tire: the physical bound on
manoeuvrability and the least accessible component to
measure~\cite{pacejka2012book}. Classical characterisation removes the vehicle
from the track---steady-state manoeuvres yield clean data at a logistical cost
a race weekend rarely affords~\cite{ontrack2025}. On-track identification
removes the logistics but reintroduces transients, noise and, critically, the
fact that a racing line is a poor excitation signal for a nonlinear model.

Dikici \emph{et al.}~\cite{ontrack2025} propose the hybrid remedy that is the
starting point of this paper. A small residual network is trained on
\SI{30}{\second} of ordinary lap data to predict the one-step error of a
nominal single-track model; the corrected model is rolled out through a
\emph{synthetic}, quasi-steady steering sweep; and Pacejka coefficients are
fitted to that sweep. Because the sweep is generated rather than driven it can
be made steady and noise-free, and because the network only has to interpolate
within its training distribution its generalisation weakness is contained. The
procedure iterates, each fit becoming the next nominal model, and on a 1:10
platform it reports a \num{3.3}$\times$ lower one-step RMSE than nonlinear
least squares under noise.

This paper asks what the pipeline identifies when it moves from a
\SI{3.7}{\kilo\gram}, \SI{0.32}{\meter}-wheelbase platform to a full-scale
Formula Student vehicle: \SI{269.6}{\kilo\gram}, \SI{1.53}{\meter} wheelbase,
\SI{16}{\degree} of road-wheel lock, on a reference profile of
\SIrange{9.6}{25.8}{\meter\per\second}. The question is identifiability rather
than parameter transfer, and two scale-dependent mechanisms, both latent at
1:10, break the estimate. First, excitation collapses inside the
virtual-data step: the Pacejka fit never sees the recorded lap, only the
synthetic rollout, whose reachable friction
$\mu_{\mathrm{reach}}\!\approx\!v^{2}\delta_{\mathrm{sweep}}/(gL)$ is a
property of the rollout configuration and not of the tire. A rollout speed inherited from the scaled platform caps it at $0.43$ here;
below the tire's peak the Magic Formula degenerates to $F_{y}\approx
F_{z}(BCD)\alpha$, only $BCD$ is identifiable, and the bounded optimiser
resolves the remaining freedom on a box edge~\cite{ljung1999sysid}. Second, lateral
demand scales as $v^{2}$, so a bad grip estimate is silent until it is
catastrophic: one identified set with a peak friction coefficient of \num{0.40}
tracked at \SI{0.16}{\meter} of cross-track error at
\SI{11.1}{\meter\per\second} and \SI{41.89}{\meter} at
\SI{27}{\meter\per\second} (Section~\ref{sec:integration}), and a scaled
platform never reaches the speed at which that error becomes observable.

A third difference is an opportunity. Because the study is conducted in the
CARLA simulator~\cite{carla2017} through a purpose-built ROS~2 interface, the
simulator's own per-wheel telemetry---slip angle, normal load, lateral
force---is available as ground truth, so an on-track identifier can be scored
against the forces the plant actually generated. Our contributions are: an identifiability analysis of the virtual-data step
(Section~\ref{subsec:excitation}) giving the bound above and an excitation-
aware rollout, verified in Section~\ref{subsec:sweep} on a grid over rollout
speed, sweep amplitude and true peak factor with a known tire; the corrections
it implies on both sides of the fit---an extrapolation-bounded sweep, a frozen
regularisation target, and a measurement chain built on the achieved road-
wheel angle, measured mass and geometry, and a friction warm start applied
only as a floor (Sections~\ref{subsec:iterative}--\ref{subsec:measurement});
admissibility gates on derived physical quantities rather than on per-
coefficient bounds (Section~\ref{sec:integration}); and a closed-loop
evaluation against the plant's own forces, in which the corrected pipeline
preserves a peak factor within \num{0.05} of the plant's effective friction
where the inherited configuration rails on its box constraint, and halves the
lateral tracking error under MPC. All results are in simulation, one run per
configuration.
